\section{Implementation}
\label{sec:implementation}

\subsection{Chunking}

\texttt{compute\_chunks(count, N)} (\texttt{include/ring\_allreduce/chunking.hpp}, \texttt{src/chunking.cpp}) splits a buffer of \texttt{count} elements into $N$ chunks as evenly as possible: with \texttt{base = count / N} and \texttt{remainder = count \% N}, the first \texttt{remainder} chunks get size \texttt{base + 1} and the rest get size \texttt{base}, with offsets the running prefix sum of sizes. This is a small, pure function with no MPI dependency, and it is the single most consequential piece of code in this repository: every send/receive offset and every step's index arithmetic in the main algorithm is downstream of it, so an off-by-one here silently corrupts every benchmark number that follows without necessarily crashing anything. It has its own unit tests (\texttt{tests/test\_chunking.cpp}) covering an even split, a split with remainder, \texttt{count < N} (which legitimately produces trailing zero-size chunks, not an error), \texttt{count == 0}, $N=1$, non-power-of-two $N$, and a count large enough to exceed a 32-bit integer, checked purely as 64-bit arithmetic since no buffer of that size needs to be allocated to validate the offset/size computation itself.

\subsection{Non-blocking communication and templated reduction}

Every step of both phases issues one \texttt{MPI\_Isend} and one \texttt{MPI\_Irecv} before a single \texttt{MPI\_Waitall(2, ...)}. This is not a style preference: Section~\ref{sec:background} already noted that a blocking send-then-receive would serialize the two halves of every step and roughly double the measured per-step time while changing nothing about how much data moves, which would make every downstream bandwidth number wrong in a way that would not be obvious from the numbers alone. Reduction during reduce-scatter cannot happen in place on receive (MPI point-to-point does not reduce), so the incoming chunk lands in a scratch buffer and is combined into the destination chunk with an explicit elementwise loop calling the reduction functor.

The reduction operator is a small functor type, not a hardcoded addition: \texttt{SumOp} (the default and the only one this project's benchmark actually measures), \texttt{MaxOp}, \texttt{MinOp}, and \texttt{ProdOp}, each a generic \texttt{operator()(T a, T b)}. The element type $T$ is resolved to an \texttt{MPI\_Datatype} through a small trait specialized for \texttt{float} and \texttt{double}; instantiating the template for an unsupported type is a compile error rather than silently picking the wrong datatype. The scratch receive buffer is reused across calls (a function-local \texttt{thread\_local} buffer that only ever grows), specifically so that a benchmark loop calling this function thousands of times per message size does not pay a heap allocation on every single call, since that allocator overhead would land exactly in the small-message latency numbers this project is trying to measure cleanly.

\subsection{Edge cases as first-class requirements}

Three cases are handled directly in the main algorithm rather than bolted on afterward:

\begin{itemize}
  \item \textbf{$N=1$}: the function returns immediately after reading the communicator's size, before any chunk computation or communication. There is nothing to combine a single rank's buffer with, so the buffer is already the answer.
  \item \textbf{\texttt{count < N}}: some chunks are legitimately size 0. A zero-byte \texttt{MPI\_Isend}/\texttt{MPI\_Irecv} is a completely legal operation, and this implementation always issues both sides of a step's handshake regardless of payload size; skipping a step's synchronization when a chunk happens to be empty would desync the ranks whose neighbor's chunk that step is \emph{not} empty. This matters in practice, not just in principle: at $N=16$ and the smallest swept message sizes (8 to 64 bytes), most chunks in the split are size 0, so this path is exercised on every real sweep, not just in a contrived unit test.
  \item \textbf{Non-power-of-two $N$}: every index in both phases is computed with plain modular arithmetic (\texttt{(rank - step + N) \% N}, normalized into $[0, N)$); nothing in the implementation assumes $N$ is a power of two, and $N \in \{3, 5, 7\}$ are included in the correctness test matrix specifically to keep that true rather than assumed.
\end{itemize}

\subsection{Correctness validation}

Every configuration is checked against two independent references, not only \texttt{MPI\_Allreduce}: the vendor collective itself, and a separately-computed \texttt{MPI\_Reduce}-to-root followed by \texttt{MPI\_Bcast}, a different pair of collectives computing the same logical answer through different machinery. Matching only \texttt{MPI\_Allreduce} would not rule out a bug in this implementation that happens to agree with Open MPI's tuned algorithm on the tested cases by coincidence. The test matrix covers $N \in \{1,2,3,4,5,7,8,16\}$, count in $\{0, 1, N-1, N, N+1, 1024\}$ plus a count large enough to exceed a 32-bit integer (checked at the chunking-arithmetic level, see above), and both \texttt{float} and \texttt{double}. Seed values are chosen deliberately (small integers for \texttt{SumOp}, powers of two for \texttt{ProdOp}) so that every comparison is an exact equality check rather than a tolerance-based one: integer-valued and power-of-two-valued operands combine identically in IEEE~754 regardless of the order in which an implementation happens to combine them, so an exact mismatch always indicates a genuine bug rather than an artifact of floating-point reassociation.
